% method: bien-doi-fourier-khong-gian
\begin{dieukienbox}
Dùng khi: bài toán Cauchy hệ số hằng trên toàn $\mathbb{R}^n$,
dữ liệu/nguồn đủ tốt để biến đổi Fourier (ví dụ $L^2$ hoặc phân bố tăng chậm).
Là phương pháp thay thế cho Kirchhoff khi muốn xem công thức dưới dạng tích chập.
\end{dieukienbox}

\begin{nhobox}
\textbf{Ý tưởng}: lấy Fourier theo $x$:
$\hat u(\xi,t)$ thỏa ODE $\hat u_{tt} + c^2|\xi|^2 \hat u = \hat f(\xi,t)$.

Giải ODE, sau đó lấy Fourier ngược để được $u$.
\end{nhobox}

\begin{meobox}
Phương pháp này ít dùng trong bài thi dạng hình học (miền ảnh hưởng),
nhưng hữu ích khi cần kiểm tra công thức hoặc dữ liệu là hàm trơn $L^2$.
\end{meobox}
